\documentclass{article}
\usepackage{graphicx} 

\title{DataMining Report}
\author{Deen Willemsen }
\date{October 2026}

\begin{document}
\maketitle

\section{Abstract}
\section{Introduction}
Online hotel reviews play an important role in helping consumers make decisions about hotels and other services. However, the increasing ability of large language models (LLMs) to generate fluent and human-like text has created new opportunities for producing deceptive reviews. Distinguishing between genuine reviews written by humans and reviews generated by an LLM is therefore an increasingly relevant classification problem.

Ignat et al. (2025) investigate this problem in their work, \textit{MAiDE-up: Multilingual Deception Detection of AI-generated Hotel Reviews}. The authors introduce the MAiDE-up dataset, which contains 10,000 real human-written hotel reviews and 10,000 AI-generated hotel reviews across ten languages. The dataset was created to investigate linguistic differences between human-written and LLM-generated hotel reviews and to evaluate methods for detecting AI-generated reviews. \cite{ignat-etal-2025-maide}

In this study, we focus on the English-language subset of the MAiDE-up dataset. We investigate whether traditional machine learning methods can distinguish between genuine and AI-generated hotel. The review texts are represented using a bag-of-words representation, after which five different classification approaches are evaluated: Multinomial Naive Bayes, logistic regression with a Lasso penalty, a single classification tree, random forests, and gradient boosting.

The aim of this study is to compare the performance of these classifiers in detecting AI-generated hotel reviews and to investigate which textual features are most useful for distinguishing between human-written and AI-generated reviews.

\section{Data}
The data used in this study comes from the MAiDE-up (Multilingual Ai-generateD fakE reviews) dataset introduced by Ignat et al. (2025). The dataset was created to study the differences between human-written and LLM-generated hotel reviews. It contains 20,000 hotel reviews, of which 10,000 are genuine human-written reviews and 10,000 are AI-generated reviews. The dataset is balanced across language, hotel location, and sentiment. The ten languages represented in the dataset are Chinese, English, French, German, Italian, Korean, Romanian, Russian, Spanish, and Turkish. \cite{ignat-etal-2025-maide}

The genuine reviews were collected from Booking.com. The authors collected reviews from hotels in ten cities: Ankara, Beijing, Berlin, Bucharest, Madrid, New Delhi, Paris, Rome, Seoul, and Washington. To obtain a balanced dataset, they identified 250 hotels in each city and collected up to 50 reviews per hotel. The reviews were also balanced with respect to sentiment. Booking.com scores from 1 to 10 were used, where scores from 1 to 6 were considered negative and scores from 7 to 10 positive. For each language, the dataset contains 500 positive and 500 negative reviews. The authors additionally performed automatic and manual quality checks on the collected reviews. \cite{ignat-etal-2025-maide}

The AI-generated reviews were created using GPT-4. The generated reviews were constructed to follow the same distributions of language, location, and sentiment as the genuine reviews. The generation process used prompts specifying the language, hotel location, and sentiment of the review. The prompts instructed the model to generate a review consisting of an ``upside'' and a ``downside'' section, together with an overall review score. The authors also used few-shot prompting and included real hotel reviews in the generation process. Quality checks were performed by native speakers, who evaluated generated reviews for readability, syntax, and style. \cite{ignat-etal-2025-maide}

For this study, we restrict the dataset to the English-language reviews, as required by the assignment. This results in 2,000 reviews: 1,000 genuine human-written reviews and 1,000 AI-generated reviews. The resulting dataset is therefore balanced with respect to the target class. The reviews contain two text fields, an ``upside'' and a ``downside'', which correspond to the two parts of the review format used by Booking.com and the generation procedure described by Ignat et al. (2025). Both fields are used as input to our classifiers.

\section{Data Pre-Processing}
The analysis was restricted to English-language reviews. This follows the scope of the dataset specified for this study, which contains 1000 genuine and 1000 AI-generated English reviews. Restricting the analysis to English also allows the same linguistic preprocessing methods, such as English stop-word removal and WordNet-based lemmatization, to be applied consistently across all reviews. 

In the data each review was split in a positive and negative parts. To get access to the full texts of the reviews we had to manually add the negative and positive parts together. Using them as separate fields was really not an option.

We devided the data into training and validation sets using a 75\%--25\% split. This follows the experimental setup specified in the assignment. Stratified sampling was used to ensure that the proportion of genuine and AI-generated reviews remained the same in both sets. A fixed random seed of 42 was used to make the split reproducible, so that the same reviews are assigned to the training and test sets each time the analysis is performed.

Before creating the bag-of-words, the texts where lower-cased. This was done to improve the data-set as e.g. the words "Hotel" and "hotel" should be treated as the same as the capitalization doesn't effect word function. The words were subsequently lemmatized using the NLTK wordnet lemmatizer. We did this because lemmatization often provides a more trainable dataset as words that have very close related meanings but differ in inflection. Making sure these words are identified the same will create more sustainable patterns that sustain better training. Stop words where hereafter removed as they don't serve a useful function

The reviews were represented using a bag-of-words representation based on word frequencies. Each review is represented by the number of times each word occurs in that review, resulting in a sparse document-term matrix. Terms occurring in fewer than two training documents were removed using a minimum document frequency of two. Such very rare terms provide limited evidence for general classification because they occur in only a single training document. Removing these sparse terms therefore reduces the dimensionality of the feature space and can reduce the influence of features that are unlikely to generalize to unseen reviews.
\section{methods/setup}
\section{Results}
\section{Conclusion}

\begin{thebibliography}{9}

\bibitem{ignat-etal-2025-maide}
Oana Ignat, Xiaomeng Xu, and Rada Mihalcea.
\newblock MAiDE-up: Multilingual Deception Detection of AI-generated Hotel Reviews.
\newblock \textit{Findings of the Association for Computational Linguistics: NAACL 2025}, pages 1636--1653, 2025.

\end{thebibliography}


\end{document}
